
%(BEGIN_QUESTION)
% Copyright 2015, Tony R. Kuphaldt, released under the Creative Commons Attribution License (v 1.0)
% This means you may do almost anything with this work of mine, so long as you give me proper credit

Les delen med tittelen ``Distribution Automation System Example'' på side 16-17 i artikkelen ``Power System Automation'' skrevet av David Dolezilek (Schweitzer Engineering Laboratories, Inc., SEL-dokument 6091), og svar på spørsmålene nedenfor:

\vskip 10pt

Diagrammet i figur 18 kalles et {\it enlinjeskjema}. Forklar hvorfor denne forenklede skjematypen er nyttig for å beskrive effektsystemer, i stedet for å tegne alle lederne i det faktiske systemet.

\vskip 10pt

Identifiser hvilke av de fem bryterne i figur 18 som er manuelt betjent, og hvilke som er automatisk betjent. Hvordan vises denne forskjellen i enlinjeskjemaet?

\vskip 10pt

Hvor er lastene i dette enlinjeskjemaet? Lastsymboler er ikke tegnet inn, men teksten gir et hint om hvor lastene finnes i systemet.

\vskip 10pt

Hvordan bidrar vernereléene i figur 18 til å beskytte utstyret i effektsystemet ved feil?

\vskip 10pt

Forklar hvorfor systemet i figur 19 er bedre med tanke på {\it forsyningssikkerhet}.

\underbar{file i01249}
%(END_QUESTION)





%(BEGIN_ANSWER)


%(END_ANSWER)





%(BEGIN_NOTES)

Et enlinjeskjema forenkler kraftsystemskjema ved å vise én leder i stedet for alle tre. Effektflyt kan da spores på samme måte som vannstrøm i et rørsystem.

\vskip 10pt

I figur 18 er bryter 1 og 3 automatisk styrt av vernerelé, mens de tre andre (2, 4 og 5) betjenes manuelt. Automatiske brytere (effektbrytere) vises med bokssymbol, mens manuelle fraskillere vises med SPST-brytersymbol.

\vskip 10pt

Lastene ligger på utgående linje 1, 2, 3 og 4, selv om de ikke er tegnet inn i enlinjeskjemaet.

\vskip 10pt

Vernereléene registrerer overstrøm og utløser bryter 1 og/eller 3 ved feil.

\vskip 10pt

I figur 19 forbedres leveringssikkerheten ved å gjøre alle brytere automatiske og koble de tilhørende vernereléene sammen, slik at de kan koordinere handlingene sine, isolere feilrammet linje automatisk og gjenopprette effekt raskest mulig til tilstøtende linjer.

%INDEX% Leseoppgave: Power System Automation (Schweitzer Engineering Labs)

%(END_NOTES)

